% Compilar: latexmk -lualatex cola_beat_by_bit.tex
\documentclass[9pt,a4paper]{extarticle}

\usepackage{fontspec}
\usepackage[brazilian]{babel}
\usepackage[a4paper,margin=11mm,top=10mm,bottom=12mm]{geometry}
\usepackage{tabularx,booktabs,array}
\usepackage{xcolor}
\usepackage{amsmath}
\usepackage{enumitem}
\usepackage{titlesec}
\usepackage{tikz}
\usepackage[american]{circuitikz}
\usepackage[hidelinks]{hyperref}
\usetikzlibrary{positioning,calc,arrows.meta,decorations.pathreplacing}

\setmainfont{TeX Gyre Heros}
\setmonofont{DejaVu Sans Mono}[Scale=0.88]

\titleformat{\section}{\normalsize\bfseries\color{blue!55!black}}{}{0pt}{}
\titlespacing*{\section}{0pt}{7pt}{3pt}
\setlength{\parindent}{0pt}
\setlength{\parskip}{3pt}
\setlist{nosep,leftmargin=12pt}
\renewcommand{\arraystretch}{1.12}
\setlength{\tabcolsep}{4pt}
\newcolumntype{Y}{>{\raggedright\arraybackslash}X}
\newcommand{\s}[1]{\texttt{#1}}          % sinal / codigo
\pagestyle{plain}

\begin{document}

{\large\bfseries Cola de montagem --- Beat by Bit (PCS3645)}\hfill
{\small T1BA5 · top \s{beat\_by\_bit} · 27/10/2026}

% ==========================================================================
\section{Guia de montagem (DE0-CV)}

Placa desligada até terminar a montagem. Tudo vai no JP1 (GPIO\_0); o JP2 fica livre.
Conferir o pino 1 do conector pela marcação da placa.

\begin{center}
\begin{tikzpicture}[
    pin/.style={circle, draw=black!35, fill=white, inner sep=0pt, minimum size=2.3mm},
    pwr/.style={pin, draw=red!60!black, fill=red!55},
    gnd/.style={pin, draw=black, fill=black!75},
    sig/.style={pin, draw=blue!60!black, fill=blue!45},
    saida/.style={pin, draw=green!45!black, fill=green!55},
    tag/.style={font=\footnotesize, inner sep=1pt},
    num/.style={font=\tiny, text=black!55},
    ld/.style={draw=black!55}
  ]
  \def\sp{0.36}
  % \conector{nome}{x}{rotulo}: 20 linhas de 2 pinos (impares a esquerda)
  \newcommand{\conector}[3]{
    \draw[rounded corners=1pt, draw=black!35, fill=black!4]
      (#2-0.2, 0.2) rectangle (#2+0.62, -19*\sp-0.2);
    \foreach \l in {0,...,19} {
      \pgfmathtruncatemacro{\pinimpar}{2*\l+1}
      \pgfmathtruncatemacro{\pinpar}{2*\l+2}
      \node[pin] (#1-\pinimpar) at (#2, -\l*\sp) {};
      \node[pin] (#1-\pinpar)   at (#2+0.42, -\l*\sp) {};
    }
    \node[font=\footnotesize\bfseries, anchor=south] at (#2+0.21, 0.25) {#3};
  }
  % rotulos: \esq{pino}{texto} a esquerda (impares), \dir{pino}{texto} a direita (pares)
  \newcommand{\esq}[2]{\draw[ld] (#1.west) -- ++(-0.3,0) node[tag, anchor=east] {#2};}
  \newcommand{\dir}[2]{\draw[ld] (#1.east) -- ++(0.3,0) node[tag, anchor=west] {#2};}

  \conector{JP1}{7.4}{JP1 · GPIO\_0}
  \node[saida] at (JP1-1) {};
  \node[sig] at (JP1-2) {};
  \node[saida] at (JP1-4) {};  \node[sig] at (JP1-6) {};  \node[saida] at (JP1-8) {};
  \node[gnd] at (JP1-12) {};  \node[gnd] at (JP1-30) {};
  \node[pwr] at (JP1-29) {};
  \node[sig] at (JP1-16) {};
  \node[sig] at (JP1-36) {};  \node[sig] at (JP1-37) {};  \node[sig] at (JP1-38) {};
  \node[sig] at (JP1-39) {};  \node[sig] at (JP1-40) {};
  \node[num, anchor=south] at (JP1-1.north) {1};
  \node[num, anchor=north] at (JP1-39.south) {39};
  \node[num, anchor=north] at (JP1-40.south) {40};

  \esq{JP1-1}{1 · \s{uart\_tx\_out} $\to$ RXD do CP2102}
  \esq{JP1-11}{11 · 5\,V da placa: \textbf{não usar}}
  \esq{JP1-29}{29 · 3,3\,V $\to$ trilho $+$: pull-ups e HC-SR04 (\textbf{nunca 5\,V})}
  \esq{JP1-37}{37 · \s{botoes[1]} · trilha 1}
  \esq{JP1-39}{39 · \s{botoes[0]} · trilha 0}
  \dir{JP1-2}{2 · \s{uart\_rx\_in} $\leftarrow$ TXD do CP2102}
  \dir{JP1-4}{4 · \s{hcsr04\_trig} $\to$ TRIG do HC-SR04}
  \dir{JP1-6}{6 · \s{hcsr04\_echo} $\leftarrow$ ECHO do HC-SR04}
  \dir{JP1-8}{8 · \s{servo\_pwm} $\to$ sinal do servo}
  \dir{JP1-12}{12 · GND $\to$ GND do CP2102}
  \dir{JP1-16}{16 · \s{reset} (apertado = reset)}
  \dir{JP1-30}{30 · GND $\to$ trilho $-$ (GND comum)}
  \dir{JP1-36}{36 · \s{botoes[4]} · START / PAUSE}
  \dir{JP1-38}{38 · \s{botoes[3]} · trilha 3}
  \dir{JP1-40}{40 · \s{botoes[2]} · trilha 2}
\end{tikzpicture}
\hspace{4mm}
\begin{circuitikz}[scale=0.8, transform shape, font=\footnotesize]
  \node[font=\footnotesize\bfseries, anchor=south] at (1.2, 4.35) {cada entrada (6×)};
  \draw (0,4) node[left] {3,3\,V (pino 29)} -- (2.4,4);
  \draw (1.2,4) to[R, l_=10\,k$\Omega$] (1.2,2) node[circ] {};
  \draw (1.2,2) -- (2.6,2) node[right, align=left] {pino do\\botão/\s{reset}};
  \draw (1.2,2) to[push button] (1.2,0) node[ground] {};
  \node[right] at (1.45,-0.25) {GND (pino 30)};
  % CP2102: os fios cruzam (TX -> RX)
  \begin{scope}[shift={(-0.6,-2.2)}]
    \node[font=\footnotesize\bfseries, anchor=south] at (2.0, -0.05) {CP2102 (fios cruzados)};
    \node[draw=black!50, fill=black!4, rounded corners=2pt, minimum width=1.3cm,
          minimum height=2.6cm] (cp) at (3.6,-1.5) {};
    \node[font=\scriptsize\bfseries, rotate=90] at (3.95,-1.5) {CP2102 · USB};
    \foreach \y/\jp/\cpn/\cor in {-0.6/pino 1 (TX)/RXD/green!45!black,
                                  -1.5/pino 2 (RX)/TXD/blue!60!black,
                                  -2.4/pino 12 (GND)/GND/black} {
      \node[font=\scriptsize, anchor=east] at (0.6,\y) {\jp};
      \node[font=\scriptsize\ttfamily, anchor=west] at (3.0,\y) {\cpn};
    }
    \draw[-{Latex[length=1.6mm]}, thick, green!45!black] (0.65,-0.6) -- (2.95,-0.6);
    \draw[{Latex[length=1.6mm]}-, thick, blue!60!black]  (0.65,-1.5) -- (2.95,-1.5);
    \draw[thick] (0.65,-2.4) -- (2.95,-2.4);
    \node[font=\scriptsize, text=black!60, align=left, anchor=north west] at (0.0,-2.7)
      {3V3, 5V, DTR, CTS: livres};
  \end{scope}
\end{circuitikz}
\end{center}

\textcolor{red!60!black}{$\bullet$} alimentação ·
$\bullet$ GND (pinos 12 e 30, ligados entre si na placa) ·
\textcolor{blue!60!black}{$\bullet$} entrada ·
\textcolor{green!45!black}{$\bullet$} saída · $\circ$ não ligar.
Ímpares à esquerda, pares à direita.

\textbf{Protoboard.} Pino 29 (3,3\,V) no trilho $+$ e pino 30 (GND) no trilho $-$.
Para cada um dos 6 botões: uma perna no trilho $-$; na outra perna, um resistor de
10\,k$\Omega$ até o trilho $+$ e um fio até o pino do JP1 da tabela. Solto, o pino
fica em 1; apertado, vai a 0. O 3,3\,V (29) é vizinho do GND (30): um fio deslocado
dá curto.

\textbf{CP2102.} Três fios, cruzados: RXD no pino 1 (TX da FPGA), TXD no pino 2 (RX da
FPGA) e GND no pino 12. O 3V3 e o 5V do CP2102 ficam livres (ele é alimentado pelo USB).
Com o jogo rodando no PC, \textbf{LEDR3 aceso} confirma a serial nos dois sentidos.

\textbf{HC-SR04.} Quatro fios: VCC no trilho $+$ (3,3\,V, \textbf{nunca 5\,V}), GND no
trilho $-$, TRIG no pino 4 e ECHO no pino 6. Alimentado a 3,3\,V, o ECHO liga direto no
GPIO, sem divisor de tensão.

\textbf{Servo (SG90).} Alimentado pela \textbf{fonte externa de 5\,V}, nunca pela placa:
fio vermelho no $+$ da fonte, marrom no $-$ da fonte e laranja (sinal) no pino 8. Ligar
o $-$ da fonte ao trilho $-$ (GND comum, senão o sinal não tem referência). O $+$ da
fonte \textbf{não} encosta no trilho $+$ de 3,3\,V nem em nenhum pino do JP1.
Na música ele marca o tempo; na vitória, levanta e abana até sair da tela.

\begin{tabularx}{\linewidth}{@{}l l l l Y@{}}
\toprule
Sinal & Pino JP1 & Pino FPGA & Liga em & Função \\
\midrule
\s{uart\_tx\_out} & 1  & N16 & RXD do CP2102          & pacotes de 22\,B a 115\,200 baud, 60 por segundo \\
\s{uart\_rx\_in}  & 2  & B16 & TXD do CP2102          & ACK do PC (0xA5 a cada 0,25\,s) acende LEDR3 \\
\s{hcsr04\_trig}  & 4  & C16 & TRIG do HC-SR04        & dispara a medição de distância \\
\s{hcsr04\_echo}  & 6  & K20 & ECHO do HC-SR04        & pulso proporcional à distância (direto, sem divisor) \\
\s{servo\_pwm}    & 8  & K22 & sinal do servo         & PWM de 50\,Hz do metrônomo (servo na fonte externa) \\
GND               & 12 & --- & GND do CP2102          & referência comum \\
\s{reset}         & 16 & T22 & botão $+$ pull-up      & apertado (0 no pino) zera o jogo para IDLE \\
3,3\,V            & 29 & --- & trilho $+$             & alimenta os 6 pull-ups e o VCC do HC-SR04 \\
GND               & 30 & --- & trilho $-$             & botões, GND do HC-SR04 e $-$ da fonte do servo \\
\s{botoes[4]}     & 36 & T20 & botão $+$ pull-up      & START / PAUSE \\
\s{botoes[1]}     & 37 & T19 & botão $+$ pull-up      & trilha 1 · música 1 (Song of Storms) \\
\s{botoes[3]}     & 38 & T18 & botão $+$ pull-up      & trilha 3 \\
\s{botoes[0]}     & 39 & T17 & botão $+$ pull-up      & trilha 0 · música 0 (Concerning Hobbits) \\
\s{botoes[2]}     & 40 & T15 & botão $+$ pull-up      & trilha 2 · música 2 (Power Rangers) \\
\bottomrule
\end{tabularx}

{\footnotesize O debounce dispara na \textbf{subida} do sinal: com pull-up, a jogada
conta quando o botão é \textbf{solto}, e novos toques ficam bloqueados por 180\,ms.
Clock: \s{clock} = CLOCK\_50 da placa (M9), sem fio.}

\clearpage
% ==========================================================================
\section{Placa: o que cada display e LED mostra}

% Display de 7 segmentos: segmentos a..g (1 = aceso)
\tikzset{
  segon/.style={line width=2pt, line cap=round, draw=red!80!black},
  segoff/.style={line width=2pt, line cap=round, draw=black!8},
}
\newcommand{\seg}[2]{\ifnum#1=1 \draw[segon] #2;\else\draw[segoff] #2;\fi}
\newcommand{\sete}[8]{%
  \begin{scope}[shift={#1}, scale=0.62, xslant=0.1]
    \draw[rounded corners=1.5pt, draw=black!35, fill=black!2] (-0.2,-0.2) rectangle (0.8,1.2);
    \seg{#2}{(0.12,1.0)--(0.48,1.0)}%
    \seg{#3}{(0.6,0.88)--(0.6,0.6)}%
    \seg{#4}{(0.6,0.4)--(0.6,0.12)}%
    \seg{#5}{(0.12,0.0)--(0.48,0.0)}%
    \seg{#6}{(0.0,0.4)--(0.0,0.12)}%
    \seg{#7}{(0.0,0.88)--(0.0,0.6)}%
    \seg{#8}{(0.12,0.5)--(0.48,0.5)}%
  \end{scope}}
\newcommand{\hexdig}[2]{\ifcase"#2 \sete{#1}1111110\or\sete{#1}0110000\or
  \sete{#1}1101101\or\sete{#1}1111001\or\sete{#1}0110011\or\sete{#1}1011011\or
  \sete{#1}1011111\or\sete{#1}1110000\or\sete{#1}1111111\or\sete{#1}1111011\or
  \sete{#1}1110111\or\sete{#1}0011111\or\sete{#1}1001110\or\sete{#1}0111101\or
  \sete{#1}1001111\or\sete{#1}1000111\fi}
\newcommand{\apagado}[1]{\sete{#1}0000000}

Da esquerda para a direita, como na placa. Exemplo do desenho: jogo em PLAY, nenhum
botão apertado e última jogada com todos os botões soltos. Chaves SW9..0 e displays
HEX5..2 não são usados.

\begin{center}
\begin{tikzpicture}[
    ledon/.style={circle, draw=red!60!black, minimum size=3.4mm, inner sep=0pt, fill=red!75},
    ledoff/.style={circle, draw=red!60!black, minimum size=3.4mm, inner sep=0pt, fill=white},
    rot/.style={font=\scriptsize, align=center, text width=1.5cm, inner sep=1pt},
    tit/.style={font=\scriptsize\bfseries, text=black!65, inner sep=1pt}
  ]
  \def\L{1.62}
  % displays: HEX5 a HEX0
  \foreach \k [count=\i from 0] in {5,4,3,2,1,0}
    \node[tit] at (5.45+\i*0.8+0.2, -0.32) {HEX\k};
  \apagado{(5.45+0*0.8,0)}
  \apagado{(5.45+1*0.8,0)}
  \apagado{(5.45+2*0.8,0)}
  \apagado{(5.45+3*0.8,0)}
  \hexdig{(5.45+4*0.8,0)}{F}
  \hexdig{(5.45+5*0.8,0)}{2}
  \node[font=\scriptsize, text=black!65, anchor=west] at (10.4, 0.62) {HEX1: botões 3..0 na última jogada};
  \node[font=\scriptsize, text=black!65, anchor=west] at (10.4, 0.25) {HEX0: estado da FSM};
  % LEDs: LEDR9 a LEDR0
  \foreach \nome/\txt/\est [count=\i from 0] in {
      LEDR9/{START\\cru}/ledon,
      LEDR8/{botão 3\\cru}/ledon,
      LEDR7/{botão 2\\cru}/ledon,
      LEDR6/{botão 1\\cru}/ledon,
      LEDR5/{botão 0\\cru}/ledon,
      LEDR4/{jogada feita\\(pulso)}/ledoff,
      LEDR3/{link serial\\ok}/ledon,
      LEDR2/{fim do chart\\(pulso)}/ledoff,
      LEDR1/{fim da\\contagem}/ledon,
      LEDR0/{START\\(pulso)}/ledoff} {
    \node[\est] at (0.4+\i*\L, -1.3) {};
    \node[tit] at (0.4+\i*\L, -0.95) {\nome};
    \node[rot, anchor=north] at (0.4+\i*\L, -1.55) {\txt};
  }
\end{tikzpicture}
\end{center}

\begin{itemize}
\item \textbf{LEDR9..5} mostram o pino cru, antes do debounce: \textbf{aceso = solto},
  apagado = apertado. É o primeiro teste da montagem: cada botão apaga só o seu LED.
  Se um LED pisca ou fica apagado sem toque, falta o pull-up ou o fio está no pino errado.
\item \textbf{LEDR4, 2 e 0} são pulsos de 1 ciclo (20\,ns): na prática não se veem.

\item \textbf{LEDR3} acende quando o PC está recebendo os pacotes e devolvendo ACK
  (apaga 1\,s depois do último). Apagado com o jogo aberto no PC: pinos 1/2 invertidos
  no CP2102, fio do pino 2 solto, GND solto ou porta serial errada no Python.
\item \textbf{LEDR1} acende ao fim da contagem regressiva e fica aceso em PLAY e PAUSE;
  apaga ao voltar para IDLE/SELECT.
\item \textbf{HEX1} guarda \s{botoes[3:0]} no instante da jogada (bit 1 = solto) e zera em
  IDLE/SELECT. Como a jogada conta ao soltar, o normal é ver F; um bit em 0 indica
  botão ainda apertado naquele instante.
\end{itemize}

\begin{tabularx}{\linewidth}{@{}c l Y Y@{}}
\toprule
HEX0 & Estado & O que acontece & Sai com \\
\midrule
0 & IDLE      & tela inicial                                & START $\to$ 6 \\
6 & SELECT    & escolha da música (Btn0/1/2)                & botão 0..3 $\to$ 1 · START $\to$ 0 \\
1 & COUNTDOWN & contagem 3, 2, 1 (180 quadros = 3\,s)        & fim da contagem $\to$ 2 \\
2 & PLAY      & jogo rodando; trilhas 0..3 nos botões 0..3   & START $\to$ 3 · fim do chart $\to$ 4 \\
3 & PAUSE     & jogo parado                                 & START $\to$ 2 · botão 0..3 $\to$ 0 \\
4 & END\_WIN  & música terminou                             & START $\to$ 0 \\
5 & END\_LOSE & não ocorre (\s{perdeu} = 0, sem limite de erros) & START $\to$ 0 \\
\bottomrule
\end{tabularx}

\textbf{Reset} (pino 16) leva a 0 a partir de qualquer estado.

\end{document}
